%--------------------BLOQUE DE FORMATO--------------------%
\usepackage[a4paper,includehead,includefoot, left=2cm,top=1.2cm,right=2cm,bottom=1.5cm,headheight=28pt]{geometry} %margenes.
%--------------------FIN DE BLOQUE DE FORMATO--------------------%

%--------------------CONFIGURACION DE IDIOMA--------------------%
\usepackage[spanish,es-tabla]{babel} % pone el idioma en español, es-tabla para Tabla en lugar de Cuadro
\usepackage[utf8]{inputenc} % Codificación de entrada, carácteres acentuados, ñ
\usepackage[T1]{fontenc} % Codificación de fuente, habilita caracteres no latinos
\usepackage{lmodern} % Fuente compatible
%--------------------FIN DE CONFIGURACION DE IDIOMA--------------------%

%-----------------PAQUETES NUESTROS-----------------------%

\usepackage[center]{caption} %permite modificar el formato de los caption para tablas/figuras
\usepackage[colorlinks=false,linkcolor=black]{hyperref} %pone vinculos en el indice y en las referencias
\usepackage{adjustbox}
\usepackage{amsmath,amsfonts,amssymb,amscd}  %para ecuaciones
\usepackage{array} %permite escribir en columnas sin ser una tabla, por ejemplo, matrices
\usepackage{bm} %permite usar simbolos matematicos en negrita
\usepackage{bookmark}
\usepackage{booktabs} %modificación de las lineas en las tablas
\usepackage{cancel} %permite tachar terminos en ecuaciones
\usepackage{chemfig} %permite realizar representaciones en 2D de estructuras químicas
\usepackage{circuitikz}
\usepackage{comment} %para comentar grandes partes de codigo, usar \begin{comment}
\usepackage{csquotes}
\usepackage{dsfont} %letras mayúsculas para conjuntos (Ej: Naturales)
\usepackage{enumerate} %permite hacer listas numeradas y por niveles
\usepackage{etoolbox} %crea condicionales, puedo tener dos PDFs distintos desde un solo archivo fuente
\usepackage{fancyhdr} % Modificar encabezados y pies de paginas.
\usepackage{float} %para el float H en tablas.
\usepackage{graphicx} %para insertar graficos e imagenes.
\usepackage{mwe} %para usar example-image como placeholder de simulaciones
\usepackage{lipsum} %introduce bloques de texto random, usado para tests
\usepackage{listings} %para ingresar codigos de MATLAB, Python, Fortran, etc
\usepackage{mathrsfs} %letras mayúsculas especiales, cursiva, gotica, etc
\usepackage{multicol}
\usepackage{multirow} %filas múltiples en tablas
\usepackage{nicefrac} %Permite escribir fracciones en el medio del texto de una forma elegante \nicefrac{a}{b}
\usepackage{parskip} %modificación de sangría
\usepackage{setspace} %permite cambiar espacio entre lineas
\usepackage{stackrel} %Permite escribir cosas arriba o abajo de otras
\usepackage{steinmetz} %mejora la forma de escribir numeros complejos
\usepackage{subfigure}
\usepackage{tabularx}
\usepackage{tikz}
\usepackage{wrapfig}
\usepackage{xcolor} %permite escribir texto en color

%-------------------- Bibliografía --------------------%
\usepackage[backend=biber,style=ieee]{biblatex} %agrega bibliografía al documento
\addbibresource{Preambulo/Bibliografia.bib}

%--------------------Titulos y enumeración de secciones--------------------%
%titulos de secciones y subsecciones.
\makeatletter
\renewcommand\section{\@startsection{section}{1}{\z@}{-3.5ex\@plus-1ex\@minus-.2ex}{2.3ex\@plus.2ex}{\normalfont\Large\bfseries}}
\renewcommand\subsection{\@startsection{subsection}{2}{\z@}{-3.25ex\@plus-1ex\@minus-.2ex}{1.5ex\@plus.2ex}{\normalfont\large\bfseries\raggedright}}
\renewcommand\subsubsection{\@startsection{subsubsection}{3}{\z@}{-3.25ex\@plus-1ex\@minus-.2ex}{1.5ex\@plus.2ex}{\normalfont\normalsize}}
\makeatother

%enumeración de secciones y subsecciones.
\renewcommand{\labelenumi}{\arabic{section}.\arabic{enumi}}
\renewcommand{\labelenumii}{\arabic{section}.\arabic{enumi}.\arabic{enumii}}
\renewcommand{\labelenumiii}{\arabic{section}.\arabic{enumi}.\arabic{enumii}.\arabic{enumiii}}
\renewcommand{\labelenumiv}{\arabic{section}.\arabic{enumi}.\arabic{enumii}.\arabic{enumiii}.\arabic{enumiv}}

\renewcommand{\thefigure}{\arabic{section}.\arabic{figure}}  %definición enumeración de figuras.

\renewcommand{\thesubfigure}{(\alph{subfigure})} %modificación enumeración de subfiguras.

\renewcommand{\thetable}{\arabic{section}.\arabic{table}} %definición enumeración de tablas.

\setcounter{tocdepth}{2} %hasta que nivel aparece en el índice
\setcounter{secnumdepth}{3} %hasta que nivel se enumera
%niveles 1:section, 2:subsection, 3:subsubsection...

%--------------------CONFIGURACION LISTING--------------------%
\usepackage[scaled=0.85]{beramono}            % load a nice TT font
\definecolor{codegreen}{rgb}{0,0.6,0}
\definecolor{codegray}{rgb}{0.5,0.5,0.5}
\definecolor{codepurple}{rgb}{0.58,0,0.82}
\definecolor{backcolour}{rgb}{0.95,0.95,0.92}

\lstdefinestyle{mystyle}{
    backgroundcolor=\color{backcolour},   
    commentstyle=\color{codegreen},
    keywordstyle=\color{magenta},
    numberstyle=\tiny\color{codegray},
    stringstyle=\color{codepurple},
    basicstyle=\ttfamily\footnotesize,
    breakatwhitespace=false,         
    breaklines=true,                 
    captionpos=b,                    
    keepspaces=true,                 
    numbers=left,                    
    numbersep=5pt,                  
    showspaces=false,                
    showstringspaces=false,
    showtabs=false,                  
    tabsize=2
}

\newcommand{\blank}{\mathord{\rule{2.5em}{0.4pt}}}

\lstset{style=mystyle}
%--------------------FIN CONFIGURACION LISTING--------------------%
